%% parts/p02/variants.tex --- Part 2.4 variants of first-order logic.
%% Phase-4 coverage: D-P3-007 (E), D-P3-008 (pointer), D-P3-009
%% (pointer plus explicit could-row stub), D-P9-002 (P), and apparatus
%% A-038/A-091.  The full free-logic card remains owned by P8 (N-b-010),
%% while the full infinitary and generalized-quantifier entries remain owned
%% by P3 (N-a-059/N-a-060).

\chapter[Variants of First-Order Logic]{Variants: Many-Sorted, Free, and Infinitary Languages}
\label{sec:p02-variants}

\noindent\textit{Reader note.}  This is a separately authored comparison
chapter following the exact inherited first- and second-order presentations.
Each variant is fixed by the sources cited in its entry; none is an alias for
an inherited calculus or a replacement for its source semantics.

The entries below separate four extensions that are often grouped under the
phrase ``nonstandard first-order syntax.''  Strict many-sorted logic admits an
exact reduction to one-sorted first-order logic; free logics require a choice
about nondenoting terms; infinitary logic changes formula length; and a
generalized quantifier brings its own semantic parameter.  Treating these as
one system would erase the distinctions on which their metatheorems depend
\cite{Manzano1996,Lindstrom1969}.

\section{Strict many-sorted first-order logic}
\label{sec:p02-many-sorted}

\begin{hbdefinition}{def:p02-many-sorted-signature}{Strict many-sorted signature and formulas}
A \emph{strict many-sorted signature} is a tuple
\[
 \Lang_{\mathrm{MS}}
   =(\mathsf{Sort},\mathcal F,\mathcal R,\operatorname{tp}),
\]
where \(\mathsf{Sort}\ne\varnothing\) is the set of sorts and
\(\operatorname{tp}\) assigns a type to every nonlogical symbol.  For the
reduction below, \(\mathsf{Sort}\), \(\mathcal F\), and \(\mathcal R\) are at
most countable; each sort has a countably infinite variable stock, and one
fixed injection of their disjoint union into the governing countable
first-order variable supply is part of the presentation.  Larger signatures
require a correspondingly larger first-order variable supply and are not in
the scope of that theorem.  Here
\(\mathsf s,\mathsf t,\mathsf s_1,\ldots,\mathsf s_n\) range over
\(\mathsf{Sort}\), \(f\) ranges over
\(\mathcal F\), and \(P\) ranges over \(\mathcal R\).  A function symbol has
a declaration
\[
 f:\mathsf s_1\times\cdots\times\mathsf s_n\longrightarrow\mathsf s,
\]
a relation symbol has a declaration
\(P:\mathsf s_1\times\cdots\times\mathsf s_n\), and a constant is a nullary
function \(c:\longrightarrow\mathsf s\).  Each sort \(\mathsf s\) has a
disjoint stock of variables \(x^{\mathsf s},y^{\mathsf s},\ldots\).

Terms are generated sort by sort: a variable or constant declared at
\(\mathsf s\) is a term of sort \(\mathsf s\), and if the arguments have sorts
\(\mathsf s_1,\ldots,\mathsf s_n\), then \(f(t_1,\ldots,t_n)\) has sort
\(\mathsf s\).  Atomic formulas are well-typed
applications \(P(t_1,\ldots,t_n)\) and equalities \(t=u\) between terms of
the same sort.  Formulas close under the classical connectives and the sorted
quantifiers \(\forall x^{\mathsf s}\) and \(\exists x^{\mathsf s}\).
Cross-sort equality is not part of this strict convention.
\index{many-sorted first-order logic}
\index{sort}
\end{hbdefinition}

The four displayed tuple fields are descriptive components of the signature,
not four additional logical primitives: their role is exhausted by the typing
clauses just given.  The strict convention is needed below because overlapping
sort domains require a different reconstruction.

\begin{hbdefinition}{def:p02-many-sorted-structure}{Strict many-sorted structure and satisfaction}
A strict \(\Lang_{\mathrm{MS}}\)-structure \(\mathcal M\) supplies a nonempty
domain \(M_{\mathsf s}\) for every \(\mathsf s\in\mathsf{Sort}\), with
\(M_{\mathsf s}\cap M_{\mathsf t}=\varnothing\) when
\(\mathsf s\ne\mathsf t\), together with interpretations
\[
 f^{\mathcal M}:M_{\mathsf s_1}\times\cdots\times M_{\mathsf s_n}
   \longrightarrow M_{\mathsf s},
 \qquad
 P^{\mathcal M}\subseteq M_{\mathsf s_1}\times\cdots\times M_{\mathsf s_n},
 \qquad
 c^{\mathcal M}\in M_{\mathsf s}
\]
matching the declarations.  An assignment sends each \(x^{\mathsf s}\) into
\(M_{\mathsf s}\).
Term evaluation and the atomic and Boolean clauses are those of
the recursive satisfaction definition in Section~\ref{sec:fol-semantics-02},
restricted to the declared sorts,
while
\[
\begin{aligned}
 \mathcal M,s\entails\forall x^{\mathsf s}\,\varphi
 &\quad\Longleftrightarrow\quad
 \mathcal M,s[x^{\mathsf s}\mapsto a]\entails\varphi
 \text{ for every }a\in M_{\mathsf s},\\
 \mathcal M,s\entails\exists x^{\mathsf s}\,\varphi
 &\quad\Longleftrightarrow\quad
 \mathcal M,s[x^{\mathsf s}\mapsto a]\entails\varphi
 \text{ for some }a\in M_{\mathsf s}.
\end{aligned}
\]
Write \(\Gamma\entails_{\mathrm{MS}}\varphi\) for consequence over all such
structures and assignments.
\index{satisfaction!many-sorted}
\end{hbdefinition}

\begin{hbdefinition}{def:p02-many-sorted-translation}{Guard translation and sort axioms}
Form a one-sorted signature \(\Lang_{\mathrm{MS}}^{\flat}\) by retaining the
nonlogical symbols without their type annotations and adjoining a unary
predicate \(Q_{\mathsf s}\) for every sort \(\mathsf s\).  Regard the disjoint sorted variable
stocks as distinct untyped variable names.  The metalinguistic map
\(\operatorname{tr}_{\mathrm{MS}}\) is the unique recursion on well-sorted
formulas that leaves atoms unchanged, commutes with the Boolean connectives,
and satisfies
\[
\begin{aligned}
 \operatorname{tr}_{\mathrm{MS}}(\forall x^{\mathsf s}\,\varphi)
   &\defeq \forall x^{\mathsf s}\bigl(Q_{\mathsf s}(x^{\mathsf s})\to
        \operatorname{tr}_{\mathrm{MS}}(\varphi)\bigr),\\
 \operatorname{tr}_{\mathrm{MS}}(\exists x^{\mathsf s}\,\varphi)
   &\defeq \exists x^{\mathsf s}\bigl(Q_{\mathsf s}(x^{\mathsf s})\land
        \operatorname{tr}_{\mathrm{MS}}(\varphi)\bigr).
\end{aligned}
\]
It is an eliminable translation, not an object-language connective.

Let \(\mathsf{Ax}_{\mathrm{sort}}(\Lang_{\mathrm{MS}})\) be the one-sorted
theory containing exactly the following sort axioms; the function-closure row
ranges over positive-arity function symbols, while constants use the preceding
row:
\[
\begin{array}{ll}
 \exists x\,Q_{\mathsf s}(x) & (\mathsf s\in\mathsf{Sort}),\\[2pt]
 Q_{\mathsf s}(c) & (c:\longrightarrow\mathsf s),\\[2pt]
 \displaystyle
 \forall x_1\cdots\forall x_n
 \left(\bigwedge_{i=1}^{n}Q_{\mathsf s_i}(x_i)
 \to Q_{\mathsf s}(f(x_1,\ldots,x_n))\right)
 & (f:\mathsf s_1\times\cdots\times\mathsf s_n\longrightarrow\mathsf s),\\[6pt]
 \forall x\,\neg\bigl(Q_{\mathsf s}(x)\land Q_{\mathsf t}(x)\bigr)
 & (\mathsf s\ne\mathsf t).
\end{array}
\]
The first family preserves the declared nonemptiness of every sort, the third
preserves function codomains, and the fourth preserves strict disjointness.
No relation-typing axiom is needed for truth of translations of well-sorted
formulas, since every translated relation atom is reached under the appropriate
guards.  One may add relation-typing axioms as an encoding convention, but
they are not members of the minimal theory defined here.

\index{guard translation!many-sorted}
\index{sort axiom}
\end{hbdefinition}

\begin{hbdefinition}{def:p02-ms-variable-guard}{Total sorted-assignment guard}
For a strict many-sorted signature \(\Lang_{\mathrm{MS}}\), let
\(Q_{\mathsf s}\) be the unary predicate of its flat translation that marks
sort \(\mathsf s\).  Because assignments are total, guarding only variables
free in selected formulas would leave every other sorted variable's value
unconstrained.  Define the metalinguistic premise set
\[
 \mathsf{VarGuard}_{\mathrm{MS}}(\Lang_{\mathrm{MS}})
 \defeq\{Q_{\mathsf s}(v):\mathsf s\in\mathsf{Sort},
       \ v\text{ is a variable of sort }\mathsf s\}.
\]
This is an assignment guard, not an object-language operator or schema.  It
differs from the bound-variable guards inserted by
\(\operatorname{tr}_{\mathrm{MS}}\) and from
\(\mathsf{Ax}_{\mathrm{sort}}\), which constrains structures rather than
assignments.  The reduction theorem below and the later matched second-order
calculus use this exact premise set.
\index{assignment guard!many-sorted}
\end{hbdefinition}

\begin{hbtheorem}{thm:p02-many-sorted-translation}{proved=annex:p2.hol.5}{Reduction of strict many-sorted consequence to one-sorted consequence}
Let \(\Lang_{\mathrm{MS}}\) be a strict many-sorted signature and let
\(\Gamma\cup\{\varphi\}\) be a set of well-sorted formulas.  Then
\[
 \Gamma\entails_{\mathrm{MS}}\varphi
 \quad\Longleftrightarrow\quad
 \mathsf{Ax}_{\mathrm{sort}}(\Lang_{\mathrm{MS}})
   \cup\mathsf{VarGuard}_{\mathrm{MS}}(\Lang_{\mathrm{MS}})
   \cup\operatorname{tr}_{\mathrm{MS}}[\Gamma]
 \entails_{\mathrm{FOL}}
   \operatorname{tr}_{\mathrm{MS}}(\varphi).
\]
The left side uses nonempty pairwise-disjoint sort domains and same-sort
equality; the right side uses ordinary nonempty-domain first-order semantics.
\index{many-sorted first-order logic!reduction}
\end{hbtheorem}

The construction behind the theorem takes the tagged union of the sort
domains, adjoins one untyped junk point, and sends every ill-typed
positive-arity function input to that point.  This is choice-free; constants
remain in their declared components.  Conversely, a model of the sort axioms is restricted to the sets
defined by its \(Q_{\mathsf s}\)'s.  The closure axioms make the restricted functions
well typed; extra one-sorted elements are ignored by the guards.  This semantic
reduction transfers compactness and downward L\"owenheim--Skolem with the
usual language bounds.  Upward enlargement transfers when a specified sort
is infinite; infinitude of the total union alone is insufficient when a
signature has infinitely many sorts.  A correspondence between two native proof
calculi would require an additional derivation translation and is not asserted
here \cite{Manzano1996}.

\section{Free-logic families and definite descriptions}
\label{sec:p02-free-logic}

Free logic concerns singular terms that may fail to denote an existing object;
it is not merely first-order logic with an empty quantifier domain.  For the
comparison in this section, fix a bivalent outer-domain presentation whenever
one is applicable: \(D=D^{\mathrm{in}}\cup D^{\mathrm{out}}\), where
\(D^{\mathrm{in}}\ne\varnothing\),
\(D^{\mathrm{in}}\cap D^{\mathrm{out}}=\varnothing\), the quantifiers range
over the inner domain \(D^{\mathrm{in}}\), and values in
\(D^{\mathrm{out}}\) represent nonexisting objects.  A distinguished
existence predicate is interpreted by
\[
 \mathcal M,s\entails E!(t)
 \quad\Longleftrightarrow\quad
 \operatorname{val}^{\mathcal M}_{s}(t)\in D^{\mathrm{in}}.
\]
This presentation does not define every free logic: partial-term systems and
neutral many-valued systems need their own semantic clauses
\cite{GratzlPavlovicNolt2025}.
\index{free logic}
\index{existence predicate}

\begin{hbtheorem}{thm:p02-free-logic-families}{cited=GratzlPavlovicNolt2025}{Three free-logic families}
Positive, negative, and neutral free logics are families distinguished first
by their treatment of atomic formulas containing a nonexisting or nondenoting
term.  Positive systems permit such an atom to be true or false; negative
systems make every ordinary such atom false; neutral systems assign a gap or
additional truth value.  These labels alone do not determine identity,
quantifier rules, connective behavior, a consequence relation, or a theory of
definite descriptions.
\index{free logic!positive}
\index{free logic!negative}
\index{free logic!neutral}
\end{hbtheorem}

In standard positive and negative outer-domain presentations, familiar
instantiation and existential-introduction formulas take guarded forms such as
\[
 \bigl(\forall x\,\varphi\land E!(t)\bigr)\to\varphi[t/x],
 \qquad
 \bigl(\varphi[t/x]\land E!(t)\bigr)\to\exists x\,\varphi,
\]
with capture avoidance and the corresponding eigenvariable restrictions in a
chosen calculus.  The guards record that quantifiers range over
\(D^{\mathrm{in}}\); they are not a common complete rule system for all three
families.

\begin{landscape}
\begingroup\scriptsize\setlength{\tabcolsep}{2pt}
\begin{longtable}{@{}p{.90in}p{1.65in}p{1.55in}p{1.55in}p{1.50in}@{}}
\caption{Free-logic variants.}
\label{tab:p02-free-logic-variants}\\
\toprule
Variant & Atom containing a\newline nonexisting term & Existence and\newline quantification & Further\newline commitment & Definite\newline descriptions \\ \midrule
\endfirsthead
\toprule
Variant & Atom containing a\newline nonexisting term & Existence and\newline quantification & Further\newline commitment & Definite\newline descriptions \\ \midrule
\endhead
Positive &
May be true or false; predicate extensions may contain outer-domain objects. &
Quantifiers range over \(D^{\mathrm{in}}\); instantiation and existential
introduction are guarded by \(E!\). &
Identity involving nonexistents may be true under the selected identity and
term semantics; there is no blanket atomic-strictness principle. &
Russellian or Lambertian principles may be added, but positivity does not force
them; in particular the contextual law below is an additional commitment. \\
\addlinespace
Negative &
Every ordinary atomic formula containing a nonexisting term is false. &
The same inner-domain quantifier discipline applies; common systems take
\(E!(t)\leftrightarrow t=t\). &
Atomic strictness has the form
\(P(t_1,\ldots,t_n)\to\bigwedge_i E!(t_i)\), with a corresponding strict
identity treatment. &
Lambert's law makes a failed description nonexisting; Russell-style truth
conditions then fit the negative treatment, subject to the exact description
axioms chosen. \\
\addlinespace
Neutral &
The atom receives a truth-value gap or an additional truth value, rather than
being classified simply as true or false. &
A quantifier range must still be declared, but the two classical guarded
formulas do not determine consequence. &
One must select a truth scheme, such as weak or strong Kleene, and a matching
proof theory. &
Description formulas inherit that gap or many-valued policy; there is no
single neutral description package before truth tables and consequence are
fixed. \\
\bottomrule
\end{longtable}
\endgroup
\end{landscape}

The comparison is a table of semantic families, not three uniquely determined
calculi.  In particular, the following two description packages must remain
distinct.  In the displays, \(\varphi(x)\) and \(\psi(x)\) have at most the
shown free variable and substitution is capture avoiding.  Russellian
contextual principles may include
\[
\begin{aligned}
 E!(\iota x\,\varphi(x))
 &\leftrightarrow
 \exists x\bigl(\varphi(x)\land
   \forall y(\varphi(y)\to y=x)\bigr),\\
 \psi(\iota x\,\varphi(x))
 &\leftrightarrow
 \exists x\bigl(\varphi(x)\land
   \forall y(\varphi(y)\to y=x)\land\psi(x)\bigr),
\end{aligned}
\]
whereas Lambert's law is
\[
 \forall x\left(
   x=\iota y\,\varphi(y)
   \leftrightarrow
   \forall y\bigl(\varphi(y)\leftrightarrow y=x\bigr)
 \right).
\]
These are theory-forming principles, not consequences of the words
``positive,'' ``negative,'' or ``neutral.''  The contextual analysis begins
with \cite{Russell1905}; the free-logic and description variants are treated
in \cite{Lambert1967,Bencivenga1986,Lehmann2002,GratzlPavlovicNolt2025}.

The broader free-logic, deontic, and conditional survey is
Section~\ref{sec:p08-conditionals-deontic-free}.  The present table remains a
semantic-family comparison rather than pretending that one family name fixes a
unique calculus.

\section{Infinitary formulas and generalized quantifiers}
\label{sec:p02-infinitary-generalized}

\begin{hbdefinition}{def:p02-lomega-one-omega}{The language \(L_{\omega_1\omega}\)}
The infinitary language \(L_{\omega_1\omega}\) extends first-order formation
as follows.  If \(\Phi\) is a countable set of formulas and
\(\bigcup_{\psi\in\Phi}\FV(\psi)\) is finite, then
\(\bigwedge\Phi\) and \(\bigvee\Phi\) are formulas, with that union as their
free-variable set.  First-order quantification is added one variable at a time,
so each quantifier string remains finite.  Satisfaction extends the
first-order clauses by
\[
 \mathcal M,s\entails\bigwedge\Phi
 \Longleftrightarrow (\forall\psi\in\Phi)\,\mathcal M,s\entails\psi,
 \qquad
 \mathcal M,s\entails\bigvee\Phi
 \Longleftrightarrow (\exists\psi\in\Phi)\,\mathcal M,s\entails\psi.
\]
The first subscript controls Boolean arity; the second controls quantifier
strings.  The finite-total-free-variable condition is part of the grammar.
\index{infinitary logic}
\index{Lomegaoneomega@\(L_{\omega_1\omega}\)}
\end{hbdefinition}

\begin{hbtheorem}{thm:p02-lomega-one-omega-calibration}{cited=Scott1965}{Basic \(L_{\omega_1\omega}\) calibration}
The language \(L_{\omega_1\omega}\) is not compact.  For every countable
structure in a countable signature there is an \(L_{\omega_1\omega}\)-sentence
whose countable models are exactly its isomorphic copies.  This Scott-sentence
claim is restricted to countable models; it is not blanket categoricity among
structures of every cardinality.
\end{hbtheorem}

For the compactness failure, if \(P_n\) is nullary for every \(n<\omega\), then
\[
 \left\{\bigvee_{n<\omega}P_n\right\}
 \cup\{\neg P_n:n<\omega\}
\]
is finitely satisfiable but not satisfiable.  Infinitary proof systems exist,
but any completeness statement must name their genuinely infinitary
derivations and semantic target; ordinary finitary strong completeness is not
restored \cite{Karp1964,Scott1965}.

\begin{hbdefinition}{def:p02-infinitary-derivation}{Infinitary derivation boundary}
For the Karp presentation used here, an infinitary derivation is a well-founded
proof tree whose nodes are labelled by \(L_{\omega_1\omega}\) sequents; a
countable-conjunction introduction may have one immediate predecessor for each
conjunct, while all quantifier strings and eigenvariable conditions remain
finite. The term is needed because an ordinary finite proof tree cannot support
that rule. It stands as a proof object, differs from a finitary derivation by
its countable branching, and is used only with the full Karp rule/axiom system
specified in the cited source \cite{Karp1964}.
\index{infinitary derivation}
\end{hbdefinition}

\begin{hbtheorem}{thm:p02-karp-completeness}{cited=Karp1964}{Karp completeness boundary}
For the Karp infinitary calculus for \(L_{\omega_1\omega}\), semantic validity
implies derivability by a well-founded infinitary proof. This is weak
completeness for the declared semantics and proof objects; it does not restore
compactness, nor does it say that arbitrary semantic consequence from an
uncountable premise set has a countable subproof.
\end{hbtheorem}

For a countable structure \(A\), enumerate its finite tuples. Define formulas
\(\theta^\alpha_{\bar a}(\bar x)\) by transfinite recursion: stage zero records
the complete atomic type of \(\bar a\); successor stages conjoin, for every
element \(b\in A\), an existential extension realized by \((\bar a,b)\), and
also assert that every possible new element realizes one of those extension
types; limit stages conjoin all earlier stages. Countability keeps each
conjunction in \(L_{\omega_1\omega}\). At the first stage where the tuple
types stabilize, the sentence for the empty tuple is a Scott sentence of
\(A\): its countable models are exactly the copies of \(A\)
\cite{Scott1965}. The restriction to countable models remains essential.

\begin{hbdefinition}{def:p02-generalized-quantifier}{Isomorphism-invariant generalized quantifiers}
For \(n\geq1\), an \(n\)-ary generalized quantifier \(Q\) is an
isomorphism-closed class of pairs \((M,R)\), where \(M\ne\varnothing\) and
\(R\subseteq M^n\).  If \(\bar x=(x_1,\ldots,x_n)\) is a tuple of pairwise
distinct variables and \(\varphi\) is a
formula, then \(Q\bar x\,\varphi\) is a formula, with satisfaction clause
\[
 \mathcal M,s\entails Q\bar x\,\varphi
 \quad\Longleftrightarrow\quad
 \left(M,
   \{\bar a\in M^n:\mathcal M,s[\bar x\mapsto\bar a]\entails\varphi\}
 \right)\in Q.
\]
For a fixed arity-indexed family \(\mathcal Q\),
\(\mathsf{FOL}(\mathcal Q)\) is first-order logic closed under exactly those
additional formation and satisfaction clauses.  The family parameter is
necessary because different quantifiers have different compactness,
L\"owenheim--Skolem, and axiomatizability behavior.  It is semantic data, not
an object-language relation symbol, and the distinction is used in the table
below \cite{Mostowski1957,Lindstrom1966,Keisler1970}.
\index{generalized quantifier}
\end{hbdefinition}

\section{Extension and hierarchy lookups}
\label{sec:p02-extension-lookup}

\begin{landscape}
\begingroup\tiny\setlength{\tabcolsep}{2pt}
\begin{longtable}{@{}p{.85in}p{1.25in}p{.80in}p{1.05in}p{1.55in}p{1.55in}@{}}
\caption{Extensions of first-order logic and surviving metatheorems.}
\label{tab:p02-fol-extensions}\\
\toprule
Extension and semantics & Added resource & Compactness & L\"owenheim--Skolem calibration & Completeness or axiomatization & Lookup conclusion \\ \midrule
\endfirsthead
\toprule
Extension and semantics & Added resource & Compactness & L\"owenheim--Skolem calibration & Completeness or axiomatization & Lookup conclusion \\ \midrule
\endhead
Strict many-sorted FOL &
Typed variables, domains, functions, and relations. &
Preserved by Theorem~\ref{thm:p02-many-sorted-translation}. &
Downward LS transfers with the ordinary language bounds; upward enlargement
requires a specified infinite sort. &
Ordinary first-order soundness and strong completeness are available on the
translated side; a native-calculus correspondence is separate. &
A disciplined presentation of first-order-level expressivity, not a
counterexample to Lindstr\"om. \\
\addlinespace
Free logic &
Terms may fail to denote existing objects, with a positive, negative, or
neutral treatment. &
System-dependent until semantics and consequence are fixed. &
System-dependent. &
System-dependent; the family name does not select one calculus. &
Choose the variant, term semantics, identity, and consequence relation before
attaching a metatheorem. \\
\addlinespace
\(L_{\omega_1\omega}\) &
Countable conjunctions and disjunctions, with finite quantifier strings. &
Fails. &
Scott-style countable-model results replace the ordinary blanket statement. &
Infinitary proof systems exist; ordinary finitary strong completeness does not. &
Greater formula length buys expressive power at the cost of compactness. \\
\addlinespace
\(\mathsf{FOL}(\mathcal Q)\) &
A specified arity-indexed family of isomorphism-invariant generalized quantifiers. &
Depends on \(\mathcal Q\). &
Depends on \(\mathcal Q\); the uncountability quantifier already blocks a blanket
downward-to-countable claim. &
Depends on \(\mathcal Q\). &
Never generalize a theorem for one quantifier to every generalized quantifier. \\
\addlinespace
Full relational second-order logic &
Relation quantifiers range over every member of \(\mathcal P(M^n)\). &
Fails. &
Both blanket downward and upward forms fail; categorical infinite
characterizations witness the obstructions. &
No sound recursively axiomatized complete calculus for full validity in an
effective language. &
Full semantics permits genuine categorical characterizations but loses the
first-order metatheorem package. \\
\addlinespace
Relational Henkin semantics with matched schemas &
Relation quantifiers range over selected nonempty families, with the stated
comprehension and choice instances. &
Preserved for an exactly matched finitary many-sorted presentation. &
Downward LS transfers with the ordinary bounds; upward enlargement requires
an infinite object sort. &
Strong completeness only for an exactly matched calculus, model class, and
translation or a separately proved native correspondence. &
The metatheorems return because the semantics is many-sorted and
model-relative; full categoricity does not. \\
\bottomrule
\end{longtable}
\endgroup
\end{landscape}

The table calibrates extensions against Lindstr\"om's theorem, but does not
apply that theorem without its hypotheses: for a regular abstract logic
extending first-order logic, closed under Boolean connectives, quantification,
and relativization, and with satisfaction invariant under isomorphism,
retaining both countable compactness and downward L\"owenheim--Skolem to
countable structures precludes a strict increase in expressive power over
first-order logic \cite{Lindstrom1969}.  The public Part~3 destination is
Theorem~\ref{thm:p03-lindstrom}.

Formulas of second-order arithmetic are classified by the
\emph{lightface analytical hierarchy}.  Its levels count alternating blocks of set or function
quantifiers over an arithmetical matrix.  For \(n\geq1\), \(\Sigma^1_n\)
begins existentially and \(\Pi^1_n\) universally, with \(n\) alternating
second-order blocks;
first-order number quantifiers inside the matrix do not raise this level.
Boldface \(\boldsymbol{\Sigma}^1_n,\boldsymbol{\Pi}^1_n\) permit real or set
parameters and must not be identified with the lightface effective classes.
The full lightface table and its basic results are in
Section~\ref{sec:p04-analytical-hierarchy} and
Table~\ref{tab:p04-analytical-hierarchy}
\cite{Kleene1952,Rogers1967,Moschovakis2009}.
\index{analytical hierarchy}
\index{lightface pointclass}
\index{boldface pointclass}
